% Libro 3 · Mundo 6 · Conectores y función de partes (app: pu-conectores)
\mundo{Conectores y función de partes}{pu-conectores}{Los conectores son las bisagras del texto: «sin embargo» anuncia un
giro, «por lo tanto» una consecuencia, «es decir» una aclaración. Leerlos bien es leer la lógica del autor. Y cada
párrafo, como cada signo de puntuación, cumple una función que el examen pregunta con frecuencia.}{%
\item \textbf{6.1} Conectores: la lógica en una palabra
\item \textbf{6.2} Función de párrafos y signos}

\leccion{pu-cn-1}{Conectores}{Juego de poner la bisagra correcta entre dos ideas.}

\begin{reto}
«El examen fue difícil; \underline{\hspace{18mm}}, casi todos aprobaron.» ¿Qué conector completa mejor la oración?
\opcionesd{por lo tanto}{además}{sin embargo}{es decir}
\end{reto}

\tecnica{La técnica: ¿qué relación hay entre las dos ideas?}
Antes de mirar los conectores de las opciones, se decide qué relación hay entre las ideas que se unen. Con la relación
clara, el conector sale solo.

\begin{center}\small
\begin{tabularx}{\linewidth}{@{}l>{\raggedright\arraybackslash}X@{}}
\toprule
\textbf{Relación} & \textbf{Conectores}\\\midrule
Adición & además, también, asimismo, incluso, por otra parte\\
Contraste & pero, sin embargo, no obstante, en cambio, aunque, a pesar de\\
Causa & porque, ya que, puesto que, debido a, dado que\\
Consecuencia & por lo tanto, por eso, así que, en consecuencia, de modo que\\
Aclaración & es decir, o sea, en otras palabras, esto es\\
Ejemplo & por ejemplo, tal como, así\\
Condición & si, siempre que, a menos que, con tal de que\\
Conclusión & en conclusión, en síntesis, finalmente, en suma\\
Orden & primero, luego, después, por último\\\bottomrule
\end{tabularx}
\end{center}

\subsubsection{Causa y consecuencia, al revés}
«Llovió, \emph{por eso} se canceló el partido» y «Se canceló el partido \emph{porque} llovió» dicen lo mismo, pero el
conector va en lugar distinto: «por eso» va antes de la consecuencia; «porque», antes de la causa.

\figura{senales}{Los conectores son las señales del camino.}

\begin{regla}
\begin{itemize}
  \item Primero la relación, luego el conector.
  \item Contraste: la segunda idea va «contra» lo que se esperaba de la primera.
  \item «Porque» introduce la causa; «por eso», la consecuencia.
\end{itemize}
\end{regla}

\begin{texto}{El valor de equivocarse}{Ensayo breve · texto original escrito para este libro}
\pa En la escuela aprendemos que el error es algo que hay que evitar. \emph{Sin embargo}, en la ciencia el error es una
herramienta: cada experimento fallido descarta una hipótesis y acerca la respuesta.

\pa \emph{Por ejemplo}, antes de dar con un filamento que funcionara, los inventores de la bombilla probaron cientos de
materiales. Ninguno de esos intentos fue un desperdicio, \emph{ya que} cada uno enseñaba qué no hacer.

\pa El problema no es equivocarse, \emph{sino} no aprender del error. Quien repite la misma equivocación sin preguntarse por
qué no está experimentando: está insistiendo. \emph{Por eso}, más que castigar el error, convendría enseñar a analizarlo.

\pa \emph{En suma}, equivocarse bien es una habilidad: la de convertir cada fracaso en información.
\end{texto}

\subsection{Preguntas comentadas}
\begin{comentada}{En el \S\,2, «ya que» introduce}
\mal{A}{una consecuencia.}{La consecuencia sería «no fueron un desperdicio»; «ya que» va antes de la razón.}
\bien{B}{la causa de que los intentos no fueran un desperdicio.}{«Cada uno enseñaba qué no hacer» es la razón.}
\mal{C}{un contraste.}{No hay oposición.}
\mal{D}{un ejemplo.}{El ejemplo lo introduce «Por ejemplo».}
\end{comentada}

\taller
\piso{Como en el examen}{opción múltiple}
\begin{ejercicios}
\item En el \S\,1, «Sin embargo» puede reemplazarse sin cambiar el sentido por
\opcionesd{No obstante}{Por lo tanto}{Además}{Es decir}
\sol{A. Los dos marcan contraste.}
\item En el \S\,3, «Por eso» introduce
\opcionesl{la causa de los errores.}{una consecuencia de lo dicho: convendría enseñar a analizar el error.}{un ejemplo.}{una contradicción.}
\sol{B.}
\item «En suma» \pf{4} anuncia
\opcionesd{una conclusión}{una excepción}{un nuevo tema}{una pregunta}
\sol{A.}
\end{ejercicios}

\begin{retoresuelto}
Lo esperable después de «difícil» sería que muchos perdieran; la segunda idea va en contra: contraste. \textbf{Opción C.}
\end{retoresuelto}
\bitacora

\leccion{pu-cn-2}{Función de párrafos y signos}{Juego de nombrar lo que hace cada parte.}

\begin{reto}
En un ensayo, un párrafo empieza: «Podría objetarse que la tecnología también aísla a las personas.» Ese párrafo
probablemente cumple la función de
\opcionesd{presentar el tema}{plantear una objeción}{dar un dato}{concluir}
\end{reto}

\tecnica{La técnica: verbos de función}
Las preguntas de función piden nombrar lo que \emph{hace} una parte del texto, no lo que dice. Sirve tener a mano una lista
de verbos: \emph{introduce, define, ejemplifica, explica, compara, contrasta, objeta, concede, refuta, concluye, resume,
cita}.

\subsubsection{Los signos también tienen función}
\begin{itemize}
  \item \textbf{Comillas}: citan palabras de otro, señalan un uso irónico o un término especial.
  \item \textbf{Paréntesis} y \textbf{rayas}: agregan una aclaración o un dato secundario.
  \item \textbf{Dos puntos}: anuncian una explicación, una enumeración o una consecuencia.
  \item \textbf{Punto y coma}: separa ideas relacionadas que podrían ser oraciones.
  \item \textbf{Signos de interrogación} en un ensayo: suelen ser preguntas retóricas que plantean el problema.
\end{itemize}

\begin{regla}
\begin{itemize}
  \item Responde «¿qué hace esta parte?» con un verbo de función.
  \item Compara la parte con lo que viene antes y después: una objeción suele ir seguida de una respuesta.
  \item Los signos aclaran (paréntesis, rayas), citan o ironizan (comillas) y anuncian (dos puntos).
\end{itemize}
\end{regla}

\begin{texto}{¿Para qué sirve la poesía?}{Ensayo · texto original escrito para este libro}
\pa ¿Para qué sirve la poesía? La pregunta suele hacerse con cierto desdén, como si la respuesta obvia fuera «para nada».
En un mundo que mide todo por su utilidad, un poema parece un lujo: no cura, no produce, no se vende.

\pa Pero la poesía cumple funciones que no figuran en ninguna estadística. Nombra lo que no tenía nombre; ordena el dolor
para que se pueda mirar; afina el oído para escuchar el idioma de otra manera. Quien lee poesía aprende a leer todo lo
demás con más atención.

\pa Podría objetarse que hoy casi nadie lee poemas. Es cierto, al menos en los libros. Pero la poesía no ha desaparecido:
se ha mudado a las canciones, a los grafitis, a las frases que alguien escribe en una servilleta y guarda durante años.

\pa La pregunta del comienzo, entonces, estaba mal planteada. La poesía no sirve \emph{para} algo, como un martillo; sirve
como sirve un mapa: para no perderse.
\end{texto}

\subsection{Preguntas comentadas}
\begin{comentada}{¿Qué función cumple el tercer párrafo?}
\mal{A}{Presenta la tesis por primera vez.}{La tesis se insinúa en el \S\,2.}
\bien{B}{Plantea una objeción y la responde.}{«Podría objetarse…», «Es cierto…», «Pero la poesía no ha desaparecido».}
\mal{C}{Da datos estadísticos.}{\emph{Dice lo contrario}: el \S\,2 dice que esas funciones no figuran en estadísticas.}
\mal{D}{Concluye el texto.}{Eso lo hace el \S\,4.}
\end{comentada}

\taller
\piso{Como en el examen}{opción múltiple}
\begin{ejercicios}
\item Las comillas en «para nada» \pf{1} señalan
\opcionesl{una cita textual de un poeta.}{la respuesta despectiva que se suele dar, que el autor no comparte.}{un título.}{una palabra extranjera.}
\sol{B.}
\item Los dos puntos no aparecen en el \S\,2, pero el punto y coma en «nombra lo que no tenía nombre; ordena el dolor…» sirve para
\opcionesl{separar una enumeración de funciones de la poesía.}{introducir una cita.}{marcar un contraste.}{cerrar el párrafo.}
\sol{A.}
\item El \S\,4 retoma la pregunta del comienzo para
\opcionesl{dejarla sin respuesta.}{concluir que estaba mal planteada y dar una respuesta con una comparación.}{cambiar de tema.}{citar a otro autor.}
\sol{B.}
\end{ejercicios}

\begin{retoresuelto}
«Podría objetarse que…» anuncia una objeción a la postura del autor, que suele responderse a continuación.
\textbf{Opción B.}
\end{retoresuelto}
\bitacora

\evaluacion
\begin{texto}{La tienda de la esquina}{Artículo de opinión · texto original escrito para este libro}
\pa Las tiendas de barrio parecían condenadas cuando llegaron los grandes supermercados. Tenían menos variedad, precios a
veces más altos y horarios irregulares. \emph{No obstante}, siguen ahí, y en muchos barrios son más numerosas que nunca.

\pa La razón no es solo económica. El tendero fía, conoce a sus clientes por el nombre y vende de a una pastilla de caldo o
de a un huevo, algo impensable en un supermercado. \emph{Es decir}, ofrece algo que no aparece en la factura: confianza.

\pa Hay quienes ven en esto una forma de atraso: «la tienda es el pasado», dicen. Pero esa opinión olvida que la tienda
cumple también una función social —es punto de encuentro, oficina de noticias, a veces hasta guardería improvisada— que
ningún almacén de cadena puede reemplazar.

\pa \emph{Por eso}, más que preguntarnos si la tienda sobrevivirá, deberíamos preguntarnos qué perderíamos si
desapareciera.
\end{texto}
\begin{ejercicios}
\item «No obstante» \pf{1} introduce
\opcionesl{una consecuencia de la llegada de los supermercados.}{un hecho que contrasta con lo que se esperaba.}{un ejemplo.}{una aclaración.}
\sol{B. (6.1)}
\item «Es decir» \pf{2} introduce
\opcionesd{una reformulación}{una causa}{una objeción}{una conclusión}
\sol{A. (6.1)}
\item El tercer párrafo cumple la función de
\opcionesl{presentar una opinión contraria y refutarla.}{dar la conclusión.}{enumerar productos.}{narrar la historia de una tienda.}
\sol{A. (6.2)}
\item Las rayas en el \S\,3 encierran
\opcionesl{una cita.}{ejemplos de la función social de la tienda.}{una opinión contraria.}{un dato numérico.}
\sol{B. (6.2)}
\item Las comillas en «la tienda es el pasado» \pf{3} indican
\opcionesl{palabras de quienes piensan distinto al autor.}{el título de un libro.}{una ironía del autor sobre sí mismo.}{una palabra técnica.}
\sol{A. (6.2)}
\end{ejercicios}
